\section{Procurement, Tenders and Contracts}


\begin{frame}{Division of Labour}
    
    \mytext{
    \begin{itemize}
        \item source of wealth (A. \noun{Smith})
        \item necessary in corporate projects
        \item internal and external
        \item role of contracts
    \end{itemize}
    }
\end{frame}

\begin{frame}{Contracts}
    
    \mytext{
    \begin{itemize}
        \item do et des
        \item[] (I give so that you give)
        \item pacta sunt servanda
        \item freedom to contract
        \item usefulness of contracts
        \item typically: objective before goal before person
        \item obligation before transaction
        \item financial and material fullfiment 
        \item control etc.
    \end{itemize}
    }
\end{frame}

   
\begin{frame}{Specifications}
    
    \mytext{
    \begin{itemize}
        \item Specify the \textbf{what,} not the \textbf{how}
        \item Leave process / sub-project planning to \textbf{contractors}
        \item Control the \textbf{results}
        \item Results are \textbf{interfaces, milestones}
        \item Contractors must be found / \textbf{procured}
        \item Objectives fixed in \textbf{contracts}
    \end{itemize}
    
    % \framebreak
    
    % For the following, we \textbf{assume}
    
    % \begin{itemize}
    %     \item Objectives are clear
    %     \item Requirements are clear
    %     \item Changes play minor role
    %     \item Process \textbf{is plannable}
    %     \item Results are \textbf{contractable}
    % \end{itemize}
    }
\end{frame}


\begin{frame}{The three constraints}
    \showmytikz{../shared/tikz/tikz-PM-three-constraints-procurement}
\end{frame}



% \begin{frame}{Focus on Systems and Interfaces}

%     Each trade or company or team or business function
    
%     ...is a system that communicates with other systems.
    
%     So we focus on the interfaces / handovers

%     \begin{itemize}
%         \item between trades
%         \item between contractors
%         \item between business functions
%         \item between teams
%         \item $\dots$
%     \end{itemize}
% \end{frame}



% \begin{frame}{System Thinking}
%     \input{tikz/tikz-system-milestones1}
% \end{frame}



% \begin{frame}{Handover / Committal / Milestone}
    
%     \input{tikz/tikz-system-milestones2}

%     \begin{itemize}
%         \item Milstones are handovers
%         \item Relevant for higher order system
%     \end{itemize}
% \end{frame}

% \begin{frame}{Planning Interfaces}
    
%     \input{tikz/tikz-system-milestones3}
    
%     From the larger perspective, 
    
%     only the interface / milestones / 
    
%     handovers / results count.
    
%     \framebreak

%     \input{tikz/tikz-system-milestones4}
    
%     To design these milestones,
    
%     think backwards: 
    
%     what does the next step need?
    
%     \framebreak

%     \input{tikz/tikz-system-milestones5}


    
%     Definition comes before the milestone,
    
%     control after the milestone.
% \end{frame}


\begin{frame}{Contracts}
    
    \mytext{
    \begin{itemize}
        \item Make clear what you want
        \item Make clear what you offer
        \item Definitions important
        \item Price? Time? Quality?
        \item Who is in charge?
        \item What events trigger what?
    \end{itemize}
    }
\end{frame}


\begingroup
    \ifdefined\setbeamercolor\setbeamercolor{background canvas}{bg=violet!20}\fi
    \begin{frame}{TASK 5 Minutes}
    
    \mytext{
        \begin{enumerate}
            \item In groups of two:
            \item make a contract!
        \end{enumerate}
    
    }
\end{frame}
\endgroup


% \input{break}


\begin{frame}{Service Phases / Project Phases}
    
    \mytext{
    of a typical architecture / engineering project
    
    \begin{itemize}
        \item requirement planning (Bedarfsplanung)
        \item projection (Entwurf)
        \item call for bids / tender (Ausschreibung)
        \item site management (Bauleitung)
    \end{itemize}
        
    }
\end{frame}

\begin{frame}{Procurement Managment}
    
    \mytext{
    \begin{description}
        \item[RFI] Request for information
        \item[RFQ] Request for quotation
        \item[RFP] Request for proposal
        \item[RFF] Request for feature
    \end{description}
    }
\end{frame}

\begin{frame}{Call for Bids, Invitation to Tender}
    
    \mytext{
    
    Auction / Invitation to Tender
    
    \begin{itemize}
        \item Supplier questionnaire
        \item Pre-qualification questionnaire: restricted tender
        \item ITT / CfB: different auction systems
        \item sealed bit / open bid
        \item first-price, second price
        \item Best: \textit{Vickrey auction}
    
    \end{itemize}
    
    }
\end{frame}

\begin{frame}{Call for Bids, Invitation to Tender}
    
    \mytext{
    
    Typical content of an ITT:
    
    \begin{enumerate}
        \item Introduction
        \item Project background
        \item Legal issues: proposed terms and conditions of contract
        \item Supplier response required
        \begin{itemize}
            \item Scope of the work
            \item Supplier profile
            \item Success criteria
            \item General contractual details
        \end{itemize}
        \item Timetable for choosing a supplier
        \item Requirements

    \end{enumerate}
    
    }
\end{frame}


\begin{frame}{Request for Proposal}
    
    \mytext{
A RfP typically includes

    \begin{itemize}
        \item Purpose
        \item Background
        \item Scope of Work
        \item Performance Standards
        \item Deliverable Quality
        \item Commercial Terms
        \item Payments, Incentives and Penalties
        \item Award Process
        \item Schedule
        \item Contacts
    \end{itemize}
    }
\end{frame}


% \begin{frame}{Information Economics - Principal-Agent-Problem }
% \begin{center}  
%     \begin{tikzpicture} 
        
%         % Add the states         
%         \node[draw]             (s) {$Principal$};         
%         \node[draw, right=of s] (r) {$Agent$ };
%                 % Connect the states with arrows         
%         \draw[every loop] 
%               (s) edge[bend right, auto=right]  node {hires} (r) 
%               (r) edge[bend right, auto=right]  node {serves} (s)   
%               (s) edge[loop left]               node {self-interest} (s)     
%               (r) edge[loop right]              node {self-interest} (r);    
%     \end{tikzpicture}
%     \mypause
%     \end{center}
    
%     \framebreak
%     \begin{minipage}{0.7\textwidth}
    
%         \begin{itemize}
%             \item The principal needs some external good or service which is beyond her own expertise.
        
%             \item The agent delivers this good or service, but has far more expertise about the solution, its making, and so forth, than the principal. 
        
%             \item There is an \textbf{information bias.} The result is \textbf{inefficiency.}
            
%             \item This results in higher prices and lower volumes.
        
%         \end{itemize}
    
%     \end{minipage}
%     \par\medskip
%     Source (example for energy efficiency): \cite{OgdenCharlesKay1985Tmom}

% \end{frame}

% \begin{frame}{Adverse selection, moral hazard}

%     Some examples: 
    
%     \medskip
    
%     \begin{minipage}{0.7\textwidth}
%         \begin{enumerate} 
%             \item the home-builder knows less than the architect less than the actual site-manager...
%             \item risk-prone seek insurance, risk-adverse don't
%             \item bad products need advertisement
%         \end{enumerate}
%     \end{minipage}
%     \par\medskip
%     Source (classic): \cite{GeorgeA.Akerlof1970TMfL}
% \end{frame}



% \begin{frame}{Four Stages of Competence}
%     \begin{tikzpicture}[scale=0.8]
%     \coordinate (A) at (-5,0) {};
%     \coordinate (B) at ( 5,0) {};
%     \coordinate (C) at (0,7) {};
%     \draw[name path=AC] (A) -- (C);
%     \draw[name path=BC] (B) -- (C);
%     \foreach \y/\A in {
%     0/Unconcious incompetence,
%     1/Conscious incompetence,
%     2/Conscious competence,
%     3/Unconscious competence} {
%         \path[name path=horiz] (A|-0,\y) -- (B|-0,\y);
%         \draw[name intersections={of=AC and horiz,by=P},
%               name intersections={of=BC and horiz,by=Q}] (P) -- (Q)
%               node[midway,above,align=center,text width=
%     \dimexpr(6em-\y em)*5\relax] {\A};
%     }
%     \end{tikzpicture}
%     \vfill
%     Source: own design / commons
% \end{frame}

% \begin{frame}{'Rumsfeld'-Diagramm}
%     \begin{tikzpicture}[scale=0.5]
%         \node at (-4,4) {1st order knowledge};
%         \node at (-4,3.3) {\emph{do we know it?}};
%             \node at (-1,2) {no};
%             \node at (-1,6) {yes};
%         \node at (4,-1.5) {2nd order knowledge};
%         \node at (4,-2.2) {\emph{do we know whether we know it?}};
%             \node at (2,-0.5) {yes};
%             \node at (6,-0.5) {no};
%         \draw (0,0) -- (0,8) -- (8,8) -- (8,0) -- (0,0);
%         \draw (4,0) -- (4,8) ; \draw (0,4) -- (8,4);
        
%         \node at (2,3) {known};
%             \node at (2,2.5) {unknowns};
            
%         \node at (2,7) {known};
%             \node at (2,6.5) {knowns};
            
%         \node at (6,3) {unknown};
%             \node at (6,2.5) {unknowns};
            
%         \node at (6,7) {unknown};
%             \node at (6,6.5) {knowns};
%     \end{tikzpicture}
%     \par\medskip
%     [own design / commons]
% \end{frame}

% \begin{frame}{Planning is Information Processing}
%     \begin{itemize}
%         \item You plan something \textbf{novel}
%         \item There are necessarily \textbf{unknowns}
%         \item You must start with \textbf{research}
%         \item Research is \textbf{asking questions}
%         \item You involve \textbf{people}
%         \item You must \textbf{communicate}
%     \end{itemize}
% \end{frame}

% \begin{frame}{Systems communicate}
%     \begin{tikzpicture}
%         \draw (0,0) rectangle (2,2) ;
%         \draw (3,2) rectangle (5,0) ;
%         \draw[->] (2,0.5) -- (3.2,0.5);
%         \draw[<-] (2,1.5) -- (3,1.5);
%         \draw (3.2,0.2) rectangle (4,1);
%     \end{tikzpicture}
%     \par
%     \cite{ParsonsTalcott1951Tss}, \cite{luhmann1984soziale}, \cite{RopohlGunter1979ESdT}
    
%     \framebreak
    
%     \begin{tikzpicture}[scale=0.5]
%         \path[fill=blue] (0,-2) -- (0,-3) -- (0.75,-3) -- (0.75,-5) -- (5.25,-5) -- (5.25,-4) -- (6,-4) -- (6,-2) -- (4,-2) -- (4,-4) -- (4.75,-4) -- (4.75,-4.5) -- (1.25,-4.5) -- (1.25,-3) -- (2,-3) -- (2,-2) -- cycle;
%         \path[draw,thick] (0,-1) -- (0,-3) -- (0.75,-3) -- (0.75,-5) -- (5.25,-5) -- (5.25,-4) -- (6,-4) -- (6,-1);
%         \path[draw,thick] (4,-1) -- (4,-4) -- (4.75,-4) -- (4.75,-4.5) -- (1.25,-4.5) -- (1.25,-3) -- (2,-3) -- (2,-1);
%     \end{tikzpicture}

% \end{frame}

\begin{frame}{Gather information!}
    
    \mytext{
    \begin{itemize}
        \item Identify the experts!
        \item Listen to the experts!
        \item (Re-)define goals / milestones \textbf{with them!}
        \item Planning \textbf{takes time!}
        \item \textbf{Planning concludes in contracts!}
        \item Milestone planning is contract planning!
    \end{itemize}
    }
\end{frame}


% \input{break}


\begin{frame}{Say you want...}
    
    \mytext{
    \begin{itemize}
        \item a house
        \item including what?
        \item the sketch?
        \item turn-key solution?
        \item the service of building?
        \item ...?
    \end{itemize}
    }
\end{frame}

\begin{frame}{Silence Procedure}
    
    \mytext{
     \begin{itemize}
         \item Qui tacet consentire videtur
         \item[] \textit{(he who remains silent seems to consent)}
         \item engl. Silence procedure
         \item tacit consent
         \item tacit acceptance procedure
         \item French: procédure d'approbation tacite
         \item dt.: necessary \textit{Kaufmännisches Bestätigungsschreiben}
     \end{itemize}

    \bigskip

    $\rightarrow$ Legal within reasonable bounds / customs
     
    }
\end{frame}
    

\begin{frame}{HOAI Service Phases}
    
    \mytext{
    
    German HOAI - `Honorarordnung für Architekten und Ingenieure' (\href{https://www.gesetze-im-internet.de/hoai_2013/}{https://www.gesetze-im-internet.de/hoai\_2013/})
    
    Official Scale of Fees for Services by Architects and Engineers
    
    \begin{enumerate}
        \item Establishing the basis of the project 
        \item Preliminary design
        \item Final design (systems and integration planning, tender documentation)
        \item Building permission application
        \item Execution drawings
        \item Preparation of contract award
        \item Assisting award process
        \item Project supervision (construction supervision)
        \item Project control and documentation 
    \end{enumerate}
    }
\end{frame}

\begin{frame}{HOAI Service Phases}
    
    \mytext{
    \begin{itemize}
        \item chargeable costs (anrechenbare Kosten)
        \item fee bands (Honorarzonen)
        \item material / quality bands definitions
    \end{itemize}
    }
\end{frame}

\begin{frame}{HOAI Service Phases}
    
    \mytext{
    
    `Leistungsphasen'
    
    \begin{enumerate}
        \item Grundlagenermittlung
        \item Vorplanung und Kostenschätzung
        \item Entwurfsplanung und Kostenberechnung 
        \item Genehnigungsplanung
        \item Ausführungsplanung
        \item Vorbereitung der Vergabe 
        \item Mitwirkung bei der Vergabe, Kostenanschlag
        \item Objektüberwachung
        \item Objektbetreuung 
    \end{enumerate}
    }
\end{frame}





\begin{frame}{Lessons Learned}
    
    \mytext{
    \begin{itemize}
        \item procurement and tenders 
        \item a cascading system
        \item multiple layers of procurement
        \item sub-projects typically other companies
        \item handovers / interfaces important
        \item contracts important
        \item controlling important
        \item damages / risks important
    \end{itemize}
    }
\end{frame}


 

\begin{frame}{Controlling/Management Control/Quality Management}
    
    \mytext{
    \begin{itemize}
        \item German Anglicism `Controling' 
        \begin{itemize}
            \item engl.~comptroller, auditor
        \end{itemize}
        \item Management Control
        \begin{itemize}
            \item Quality / Scope Control (`hardware')
            \item Quality assurance / inspection
            \item Financial Control / auditing ('bills')
        \end{itemize}\end{itemize}
    }
\end{frame}

\begin{frame}{Controlling/Management Control/Quality Management}
    
    \mytext{
\begin{itemize}

        \item audit \emph{work} of contractor
        \item audit \emph{contractor}
        \item ISO 9000
        \item Testing, software testing
    \end{itemize}
    }
\end{frame}


\begin{frame}{Control}
    
    \mytext{
    German `Controlling' $\neq$ engl.~auditing, control
    
    \begin{itemize}
        \item approval
        \item quality control
        \item quantity control
        \item time control
        \item financial control
        \item $\dots$
    \end{itemize}
    Control takes time!
    
    Cash flow issues!
    }
\end{frame}

\begin{frame}{Control}
    
    \mytext{
    
    \begin{description}
        \item[characteristic] to be checked
        \item[sensor] or means to control
        \item[comparator] or standard
        \item[activator] or corrective measure
    \end{description}
    }
\end{frame}



\begin{frame}{Contracts}
    
    \mytext{
    % \begin{itemize}
    %     \item do ut des
    %     \item freedom of contract (`Vertragsfreiheit', `caveat emptor')
    %     \item pacta sunt servanda
    %     \item but: legal restraints
    %     \item incentives, moral hazard
    %     \item insurance
    %     \item no contract to the detriment of a third party
    %     \item \textbf{model contracts}
    % \end{itemize}
    % \framebreak
    \begin{itemize}
        \item When does the object change ownership?
        \item When are payments happening?
        \item What are the conditions to do what?
        \item How is quality assured?
        \begin{itemize}
            \item quality, quantity, properties
            \item price, payment
            \item communications
        \end{itemize}
        \item What happens when contract is breached?
        \item Culpa in contrahendo (civil law)
        \item Estoppel (common law)
    \end{itemize}
    }
\end{frame}

\begin{frame}{Legal Systems}
    \showmybitmap{../shared/img/legalsystems-map.png}
    \mycomment{
    
    
    Source: \href{http://chartsbin.com/view/aq2}{http://chartsbin.com/view/aq2}
    }
\end{frame}

\begin{frame}{Damages and Compensations}
    
    \mytext{
Some legal terms

    \begin{itemize}
        \item damages from breaches of the contract \emph{etc.}
        \item compensatory damages (`Schadenersatz')        
        \item expectation damages
        \item reliance damages
        \item restitution (`Rückerstattung')
        \item tort (ca.~`unerlaubte Handlung') (e.g. negligence, nuisance)
        \item \emph{Nebenschäden, Folgeschäden, Spezielle Schäden, Ausfallschäden...}
        \end{itemize}
    }
\end{frame}

\begin{frame}{Damages and Compensations}
    
    \mytext{
\begin{itemize}
            
        \item contract clauses
        \item fix compensation for damages:
        \item \textbf{liquidated damages} (`pauschalierter Schadenersatz')
        \item no need to proof etc.
        \item \textbf{penal} damage illegal \emph{in common law}
                \item $\neq$ dt. `Vertragspönale, Vertragsstrafe'
    \end{itemize}
    
    }
\end{frame}



\begingroup
    \ifdefined\setbeamercolor\setbeamercolor{background canvas}{bg=violet!20}\fi
    \begin{frame}{TASK (voluntary)}
    
    \mytext{
        \begin{enumerate}
            \item Write a Request for Proposal / Call for Bids
            \item Related to your own project
            \item Research model contract(s)
        \end{enumerate}
    
    }
\end{frame}
\endgroup